% Part: normal-modal-logic
% Chapter: frame-correspondence
% Section: properties-accessibility

\documentclass[../../../include/open-logic-section]{subfiles}

\begin{document}

\olfileid{nml}{frd}{acc}

\olsection{ప్రాప్యత సంబంధాల ధర్మాలు}

అనేక మోడల్ \tetoken{సూత్రాలు}{!!{formula}s} ప్రాప్యత సంబంధానికి
ఉండే సరళమైన, కొన్నిసార్లు పరిచితమైన ధర్మాలకు లక్షణాత్మకంగా
ఉంటాయి. ఒక దిశలో దీని అర్థం: నిర్దిష్ట ధర్మం గల ఏ నమూనాలోనైనా,
దానికి అనురూపమైన \tetoken{సూత్రం}{!!{formula}}, దాని అన్ని
ప్రతిస్థాపన నిదర్శనాలూ సత్యమవుతాయి. ప్రాప్యత సంబంధాల ఐదు
సాంప్రదాయిక రకాలను, అవి సత్యమయ్యేలా హామీ ఇచ్చే సూత్రాలను
మొదట చూద్దాం.

\begin{thm}\ollabel{thm:soundschemas}
  $\mModel{M} = \tuple{W, R, V}$ ఒక నమూనా అనుకుందాం.
  $R$కు \olref{tab:five} ఎడమవైపు చూపిన ధర్మం ఉంటే,
  కుడివైపు ఉన్న \tetoken{సూత్రం}{!!{formula}} యొక్క ప్రతి
  నిదర్శనం~$\mModel{M}$లో సత్యమవుతుంది.
\end{thm}

\begin{table}[t]
    \begin{tabular}{| l || l |}
      \hline
      {\emph{$R$ ఇలా ఉంటే \dots}} & {\emph{$\mModel{M}$లో ఇది సత్యం:}} \\
      \hline \hline
      \emph{సీరియల్}: $\forall u \exists v Ruv$ & \hbox to.43\textwidth
           {$\Box p \lif \Diamond p$ \hfill (\Ax{D})} \\
      \hline
      \emph{స్వావర్తన}: $\forall w Rww$
      & \hbox to.43\textwidth
          {$\Box p \lif p$ \hfill (\Ax{T})} \\
      \hline
      \emph{సౌష్ఠవ}: &  \hbox to .43\textwidth
           {$p \lif \Box\Diamond p$ \hfill (\Ax{B})} \\
           $\forall u\forall v(Ruv \lif Rvu)$ & \\
      \hline
      \emph{సంక్రామక}: & \hbox to.43\textwidth
           {$\Box p \lif \Box \Box p$ \hfill (\Ax{4})} \\
      $\forall u \forall v \forall w ((Ruv \land Rvw) \lif Ruw)$ & \\
      \hline
      \emph{యూక్లిడియన్}: & \hbox to .43\textwidth
        {$\Diamond p \lif \Box\Diamond p$ \hfill (\Ax{5})} \\
      $\forall w \forall u \forall v ((Rwu \land Rwv) \lif Ruv)$ &\\
      \hline
    \end{tabular}
    \caption{ఐదు అనురూపతా వాస్తవాలు.}
    \ollabel{tab:five}
  \end{table}


\begin{proof}
  \Ax{B} సందర్భాన్ని చూద్దాం: ఒక నమూనాలో ఈ పథకం
  సత్యమని చూపాలంటే, దాని అన్ని నిదర్శనాలూ ఆ నమూనాలోని
  అన్ని లోకాల్లో సత్యమని చూపాలి. అందువల్ల \Ax{B}కు
  $!A \lif \Box\Diamond !A$ అనే నిదర్శనాన్ని తీసుకుందాం;
  $w \in W$ను ఏకపక్ష లోకంగా తీసుకుందాం.
  $w$ వద్ద $\Box \Diamond !A$ సత్యమని చూపడానికి,
  పూర్వపక్షం $!A$ అనేది $w$ వద్ద సత్యమని అనుకుందాం.
  కాబట్టి $w$ నుంచి ప్రాప్యమయ్యే ప్రతి $w'$ వద్ద
  $\Diamond !A$ సత్యమని చూపాలి. ఇప్పుడు $Rww'$
  అయ్యే ఏ $w'$కైనా, సౌష్ఠవ ఊహ వల్ల $Rw'w$ కూడా
  ఉంటుంది (\olref{fig:Bsymm} చూడండి). $\mSat{M}{!A}[w]$
  కాబట్టి $\mSat{M}{\Diamond !A}[w']$.
  $Rww'$ అయ్యే ఏకపక్ష లోకం $w'$ను తీసుకున్నాం కాబట్టి
  $\mSat{M}{\Box\Diamond!A}[w]$.

  మిగిలిన సందర్భాలను అభ్యాసాలుగా వదిలేస్తున్నాం.
\end{proof}

\begin{prob}
  \olref[nml][frd][acc]{thm:soundschemas} నిరూపణను పూర్తి చేయండి.
\end{prob}

\begin{figure}
  \begin{center}
    \begin{tikzpicture}[modal]
      \node[world] (w1) [label={below:$\mSat{{}}{\formula{A}}$\\
          $\mSat{{}}{\Box\Diamond \formula{A}}$}] {$w$};
      \node[world] (w2) [label=below:{$\mSat{{}}{\Diamond \formula{A}}$},
        right=of w1] {$w'$};
      \draw[->, bend left] (w1) to (w2);
      \draw[->, bend left] (w2) to (w1);
    \end{tikzpicture}
  \end{center}
\caption{సౌష్ఠవాన్ని ఉపయోగించిన వాదన.}
\ollabel{fig:Bsymm}
\end{figure}

\olref{thm:soundschemas}లోని విలోమ సూచనలు నిలవవని గమనించండి:
ఒక నమూనాలో పథకం సత్యమైతే, ఆ నమూనా ప్రాప్యత సంబంధానికి
దానికి అనురూపమైన ధర్మం తప్పక ఉంటుందని చెప్పలేం.
\Ax{T}, స్వావర్తన నమూనాల సందర్భంలో, \Ax{T} పథకానికే
ఒక నిదర్శనం అసత్యమయ్యే నమూనాను సులభంగా ఇవ్వవచ్చు:
$W = \{w\}$, $V(p) = \emptyset$ అనుకుందాం;
ప్రాప్యత సంబంధంలో ఏ జతా లేకుండా తీసుకుందాం.
\sourcecorrection{OLTENMLFRDACC-001}{మూలం ఒకే లోకాన్ని,
  $V(p)=\emptyset$ను ఇచ్చి $R$ స్వావర్తనం కాదని,
  $\Box p$ సత్యమని నిర్ధారిస్తుంది; కానీ $R$ను నిర్దేశించదు.
  ఆ నిర్ధారణలకు అవసరమైన $R=\emptyset$ను ఉదాహరణలో
  స్పష్టంగా చేర్చాం.}
అప్పుడు $R$ స్వావర్తనం కాదు; అయితే $\mSat{M}{\Box
  p}[w]$ మరియు $\mSat/{M}{p}[w]$. కానీ ఇక్కడ
\Ax{T}కు చెందిన ఒక్క నిదర్శనం మాత్రమే~$\mModel{M}$లో
అసత్యమైంది; $\Box \lnot p \lif \lnot p$ వంటి ఇతర
నిదర్శనాలు సత్యమే. \Ax{T} యొక్క \emph{ప్రతి ప్రతిస్థాపన
నిదర్శనం}~$\mModel{M}$లో సత్యమైనా $\mModel{M}$
స్వావర్తనం కాని ఉదాహరణలు ఇవ్వడం మరింత కష్టం.
అయితే అలాంటి నమూనాలు కూడా ఉన్నాయి:

\begin{prop}\ollabel{prop:reflexive}
  $\mModel{M} = \tuple{W, R, V}$ అనే నమూనాలో $W = \{u, v
  \}$ అనుకుందాం. $u$, $v$ లోకాలు $R$ ద్వారా
  ఒకదానికొకటి సంబంధించబడి ఉండాలి: అంటే $Ruv$,
  $Rvu$ రెండూ ఉండాలి. ఈ రెండు జతలే సంబంధంలో ఉండాలని
  కూడా తీసుకుందాం.
  \sourcecorrection{OLTENMLFRDACC-002}{మూలం $Ruv$,
    $Rvu$లను పేర్కొన్నా $Ruu$, $Rvv$ లేవని చెప్పదు;
    అందువల్ల చివరలో చెప్పిన అస్వావర్తనత్వం ఆ ఊహలతో
    అనుగమించదు. ఉద్దేశించిన ఉదాహరణకు
    $R=\{(u,v),(v,u)\}$ అని స్పష్టం చేశాం.}
  అన్ని $p$లకు $u \in V(p) \Leftrightarrow v
  \in V(p)$ అని అనుకుందాం. అప్పుడు:
  \begin{enumerate}
  \item అన్ని $!A$లకు: $\mSat{M}{!A}[u]$ అప్పుడే, అప్పుడే
    $\mSat{M}{!A}[v]$ ($!A$పై ఆగమనాన్ని వాడండి).
  \item \Ax{T} యొక్క ప్రతి నిదర్శనం $\mModel{M}$లో సత్యం.
  \end{enumerate}
  $\mModel{M}$ స్వావర్తనం కాదు (నిజానికి అది
  \emph{అస్వావర్తనం}) కాబట్టి, \Ax{T} సందర్భంలో
  \olref{thm:soundschemas} విలోమం విఫలమవుతుంది
  (\olref{thm:soundschemas}లో పేర్కొన్న ఇతర పథకాలలో
  కొన్ని పథకాలకూ---అన్నింటికీ కాదు---ఇలాంటి వాదనలు ఇవ్వవచ్చు).
\end{prop}

\begin{prob}
  \olref[nml][frd][acc]{prop:reflexive}లోని వాదనలను నిరూపించండి.
\end{prob}

\Ax{D}, \Ax{T}, \Ax{B}, \Ax{4}, \Ax{5} అనే ఐదు
సాంప్రదాయిక \tetoken{సూత్రాల}{!!{formula}s}పై ప్రధానంగా
దృష్టి పెట్టబోతున్నా, ప్రాప్యత సంబంధాల మరికొన్ని
ధర్మాలను \olref{tab:anotherfive}లో నమోదు చేస్తున్నాం.
ఏ లోకం నుంచైనా గరిష్ఠంగా ఒకే లోకం ప్రాప్యమైతే,
ప్రాప్యత సంబంధం~$R$ \emph{పాక్షిక ప్రమేయాత్మకమైనది}.
ఏ లోకం నుంచైనా ఖచ్చితంగా ఒకే లోకం ప్రాప్యమైతే,
దానిని \emph{ప్రమేయాత్మకమైనది} అంటాం.
(అందువల్ల ప్రమేయాత్మక సంబంధాలు సీరియల్, పాక్షిక
ప్రమేయాత్మక అనే రెండు ధర్మాలూ గలవే.) ప్రాప్యత సంబంధం
ఒక (పాక్షిక) ప్రమేయంలా పనిచేస్తుంది కాబట్టి వాటికి
``ప్రమేయాత్మక'' అనే పేరు వచ్చింది. $Ruv$
అయినప్పుడల్లా $u$, $v$ల ``మధ్య'' ఒక~$w$
ఉంటే ఆ సంబంధం \emph{బలహీన సాంద్రత గలది}.
ఈ అర్థంలో బలహీన సాంద్ర సంబంధాలు సంక్రామక సంబంధాలకు
వ్యతిరేక దిశలో ఉంటాయి: సంక్రామక సంబంధంలో $u$ నుంచి
$w$ మీదుగా $v$కు చేరగలిగితే $u$ నుంచి $v$కు
నేరుగానూ చేరగలం; బలహీన సాంద్ర సంబంధంలో $u$ నుంచి
$v$కు నేరుగా చేరగలిగితే ఏదో~$w$ మీదుగా కూడా
చేరగలం. ఏదో లోకం~$w$ నుంచి $u$, $v$ లోకాలు
రెండింటికీ చేరగలిగినప్పుడల్లా, ఆ $u$, $v$
రెండింటి నుంచీ ఒకే లోకం~$t$కు చేరగలిగితే సంబంధం
\emph{బలహీన సహగమ్యమైనది}. దీన్నే కొన్నిసార్లు
``వజ్ర ధర్మం'' లేదా ``సంగమం'' అంటారు.

\begin{table}[t]
    \begin{tabular}{| p{.50\textwidth} || p{.43\textwidth} |}
      \hline
      {\emph{$R$ ఇలా ఉంటే \dots}} & {\emph{$\mModel{M}$లో ఇది సత్యం:}} \\
      \hline \hline
      \emph{పాక్షిక ప్రమేయాత్మక}: &
      \multirow{2}{*}{$\Diamond p \lif \Box p$} \\
      $ \forall w \forall u \forall v ((Rwu \land Rwv) \lif u =v )$ & \\
      \hline
      \multirow{2}{*}{\emph{ప్రమేయాత్మక}: $\forall w \exists v \forall u(Rwu \liff \eq[u][v]) $}
      & \multirow{2}{*}{$\Diamond p \liff \Box p$} \\
      & \\
      \hline
      \emph{బలహీన సాంద్ర}: &  \multirow{2}{*}{$\Box \Box p \lif
        \Box p$} \\
      $\forall u\forall v(Ruv \lif \exists w(Ruw \land Rwv))$ & \\
      \hline
      \emph{బలహీన సంయుక్త}: &
      \multirow{3}{*}{\hbox to.43\textwidth{$
        \begin{array}{@{}l@{}}
          \Box ((p \land \Box p) \lif q) \lor {} \\
          \qquad \Box ((q \land \Box q) \lif p)
        \end{array}$ \hfill (\Ax{L})}
      } \\
      $\forall w \forall u \forall v ((Rwu \land Rwv) \lif {}$ &\\
      \qquad $(Ruv \lor u=v \lor Rvu))$ & \\
      \hline
      \emph{బలహీన సహగమ్య}: & \multirow{3}{*}{\hbox to .43\textwidth
        {$\Diamond\Box p \lif \Box\Diamond p$\hfill (\Ax{G})}} \\
      $\forall w \forall u \forall v ((Rwu \land Rwv) \lif {}$ & \\
      \qquad $\exists t (Rut \land Rvt))$ & \\
      \hline
    \end{tabular}
    \caption{ఇంకా ఐదు అనురూపతా వాస్తవాలు.}
    \ollabel{tab:anotherfive}
  \end{table}

\begin{prob}
  $\mModel{M} = \tuple{W, R, V}$ ఒక నమూనా అనుకుందాం.
  $R$కు \olref[nml][frd][acc]{tab:anotherfive} ఎడమవైపు
  చూపిన ధర్మాలు ఉంటే, అనురూపమైన కుడివైపు సూత్రం యొక్క
  ప్రతి నిదర్శనం~$\mModel{M}$లో సత్యమని చూపండి.
\end{prob}



\end{document}
